\documentclass{article}
\usepackage[margin=1in]{geometry} % margins
\usepackage[justification=centering]{caption}
%\usepackage{calrsfs}
\usepackage{amsmath}
\usepackage{amsfonts}
\usepackage{amssymb}
\usepackage{parskip}
\usepackage{graphicx}
\usepackage{enumerate}
\usepackage{enumitem}
\usepackage{amsthm}
\usepackage{float}
\usepackage{tikz-cd}
\usepackage{ytableau}
\usepackage{dsfont}

\renewcommand{\implies}{\Rightarrow}
\newcommand{\longimplies}{\Longrightarrow}
\newcommand{\N}{\mathbb{N}}
\newcommand{\Z}{\mathbb{Z}}
\newcommand{\Q}{\mathbb{Q}}
\newcommand{\R}{\mathbb{R}}
\newcommand{\C}{\mathbb{C}}
\newcommand{\D}{\mathbb{D}}
\renewcommand{\H}{\mathbb{H}}
\newcommand{\F}{\mathbb{F}}
\newcommand{\K}{\mathbb{K}}
\newcommand{\Lf}{\mathbb{L}}
\newcommand{\Sp}{\mathbb{S}}
\newcommand{\cl}[1]{\overline{#1}}
\renewcommand{\emptyset}{\varnothing}
\newcommand{\norm}[1]{\left\|#1\right\|}
\newcommand{\maxnorm}[1]{\left\|#1\right\|_{\text{max}}}
\newcommand{\opnorm}[1]{\left\|#1\right\|_{\mathrm{op}}}
\newcommand{\supnorm}[1]{\left\|#1\right\|_{\infty}}
\newcommand{\id}{\operatorname{id}}
\newcommand{\sspan}{\operatorname{span}}
\newcommand{\tr}{\operatorname{tr}}
\newcommand{\ch}{\operatorname{char}}
\newcommand{\poly}{\mathrm{P}}
\newcommand{\mtx}{\mathcal{M}}
\newcommand{\seq}{\mathbf{Seq}}
\newcommand{\vv}[1]{\mathbf{#1}}
\newcommand{\dirsum}{\oplus}
\newcommand{\rank}{\operatorname{rank}}
\newcommand{\im}{\operatorname{im}}

\newcommand{\cj}[1]{\overline{#1}}
\newcommand{\inprod}[2]{\left\langle#1,#2\right\rangle}
\renewcommand{\Re}{\operatorname{Re}}
\renewcommand{\Im}{\operatorname{Im}}

\newcommand{\res}{\operatorname{res}}

\newcommand{\floor}[1]{\left\lfloor#1\right\rfloor}

% Theorem environments
\usepackage{amsthm}

\theoremstyle{definition} % The "plain" style italicizes all body text.
	\newtheorem{thm}{Theorem}
		\numberwithin{thm}{section} % Theorem numbers are determined by section.
	\newtheorem{lemma}[thm]{Lemma}
	\newtheorem{prop}[thm]{Proposition}
	\newtheorem{cor}[thm]{Corollary}
	\newtheorem{conj}[thm]{Conjecture}

\theoremstyle{definition} % The "definition" style does not italicize body text..
    \newtheorem{defn}[thm]{Definition}
	\newtheorem{example}[thm]{Example}
    \newtheorem{exercise}[thm]{Exercise}
    \newtheorem{question}[thm]{Question}

\usepackage{mathtools}

\usetikzlibrary{decorations.markings,positioning,decorations.text}

\def\xr{5}
\def\yr{5}

%\usepackage{titlesec}
%\titleformat{\section}[hang]{\normalfont\fontsize{14}{0}\bfseries}{Problem 8.\thesection\ }{0pt}{}

\begin{document}

\section{Worksheet}

\begin{question}
Let $f : E \to [0, \infty]$ satisfy $\int_E f < \infty$. Let $A_n \subseteq A_{n-1} \subseteq E$ be a decreasing sequence of sets such that $m(A_n) \to 0$ as $n \to \infty$. Show that
\[
\lim_{n \to \infty} \int_E f \mathds{1}_{A_n} = 0
\]
\end{question}
\begin{proof}
Let $\varepsilon > 0$ and let $\varphi: E \to [0, \infty]$ be a simple function satisfying $\varphi \leq f$ and $(\int_E f) - \varepsilon < \int_E \varphi$. We then have the following chain of inequalities for any $n$:
\[ 0 \leq \int_E (f - \varphi)\mathds{1}_{A_n} \leq \int_E (f - \varphi)\mathds{1}_{A_n} + \int_E (f - \varphi)\mathds{1}_{A_n^c} = \int_E (f - \varphi) < \varepsilon \]
Now we write $\varphi = \sum_{j = 1}^k \alpha_j \mathds{1}_{E_j}$ for some $E_j \subseteq E$ and get 
\[ \int_E \varphi\mathds{1}_{A_n} = \int_E \sum_{j = 1}^k \alpha_j \mathds{1}_{E_j \cap A_n} \leq \sum_{j = 1}^k \alpha_j m(A_n) \]
which approaches 0 as $n \to \infty$. We then have 
\[ \int_E f\mathds{1}_{A_n} < \int_E \varphi\mathds{1}_{A_n} + \varepsilon \longimplies 0 \leq \limsup_{n \to \infty} \int_E f\mathds{1}_{A_n} \leq \lim_{n \to \infty}\int_E \varphi\mathds{1}_{A_n} + \varepsilon = \varepsilon \]
Since the above holds for all $\varepsilon > 0$, the limsup, hence the limit itself, of $\int_E f\mathds{1}_{A_n}$ as $n \to \infty$ is 0.
\end{proof}

\begin{question}
Let $E$ be Lebesgue measurable with $m(E) > 0$. Show that there exists $p \in E$ so that $m(E \cap \{x : |x - p| < \varepsilon \}) > 0$ for all $\varepsilon > 0$. (Problem suggested by a student.)
\end{question}
\begin{proof}
Let $\varepsilon > 0$ and consider the sets $E_n := (n\varepsilon, (n + 1)\varepsilon) \cap E$ and note that the union of all $E_n$ for $n \in \Z$ is $E \setminus \{n\varepsilon: n \in \Z\}$ which has the same measure as $E$. We then have 
\[ m(E) = \sum_{n \in \Z} m(E_n) \]
and since $m(E) > 0$, we cannot have $m(E_n) = 0$ for all $n$. Thus for some $n \in \Z$, $m(E_n) > 0$, so we can take any $p \in E_n = (n\varepsilon, (n + 1)\varepsilon) \cap E$ and we will have $(p - \varepsilon, p + \varepsilon) \supseteq (n\varepsilon, (n + 1)\varepsilon)$, hence $m(E \cap (p - \varepsilon, p + \varepsilon)) \geq m(E_n) > 0$ as required.
\end{proof}

\begin{question}
Let $f_n : \mathbb{R} \to [0, \infty]$ be a sequence of Lebesgue measurable functions that converge pointwise to some $f$. Assume that
\[
\lim_{n \to \infty} \int_{\mathbb{R}} f_n = \int_{\mathbb{R}} f < \infty
\]
Show that
\[
\lim_{n \to \infty} \int_{\mathbb{R}} \mathds{1}_E f_n = \int_{\mathbb{R}} \mathds{1}_E f
\]
for every measurable $E$. Show that this may fail if $\lim_{n \to \infty} \int_{\mathbb{R}} f_n = \int_{\mathbb{R}} f = \infty$.
\end{question}
\begin{proof}
Note that if we show that $\int_\R |f_n - f| \to 0$, then we will have 
\[ 0 \leq \left|\int_\R \mathds{1}_E f_n - \int_\R \mathds{1}_E f\right| \leq \int_\R \mathds{1}_E \left|f_n - f\right| \leq \int_\R \left|f_n - f\right| \to 0 \]
and hence the desired limit identity holds. For $\int_\R |f_n - f| \to 0$, we note that $f + f_n - |f_n - f|$ is either $f$ or $f_n$ for each $x$, hence nonnegative, and its liminf is $2f$. We thus have by Fatou's lemma
\begin{align*}
    2\int_\R f 
    &= \int_\R \liminf_{n \to \infty} (f + f_n - |f_n - f|) \\
    &\leq \liminf_{n \to \infty} \int_\R (f + f_n - |f_n - f|) \\
    &= \int_\R f + \liminf_{n \to \infty} \int_\R f_n - \limsup_{n \to \infty} \int_\R |f_n - f| \\
    &= 2\int_\R f - \limsup_{n \to \infty} \int_\R |f_n - f| \\
    \longimplies 0 \leq \lim_{n \to \infty} \int_\R |f_n - f| &\leq \limsup_{n \to \infty} \int_\R |f_n - f| \leq 0
\end{align*}
Thus we see that $\lim_{n \to \infty} \int_\R |f_n - f| = 0$, which completes the proof.
\end{proof}

\begin{question}
Let $f_n : E \to [0, \infty]$ be a decreasing sequence of Lebesgue measurable functions such that $\int_E f_1 < \infty$. Let $f$ be the pointwise limit. Show that
\[
\lim_{n \to \infty} \int_E f_n = \int_E f
\]
\end{question}
\begin{proof}
Just apply DCT with $g = f_1$.
\end{proof}

\begin{question}
If $f : \mathbb{R} \to \overline{\mathbb{R}}$ is integrable, show that the function
\[
F(x) := \int_{\mathbb{R}} \mathds{1}_{(-\infty, x)} f
\]
is continuous.
\end{question}
\begin{proof}
Let $x \in \R, \varepsilon > 0$, and consider the measurable sets $A_n = [x, x + 1/n)$, which is a decreasing sequence with $m(A_n) \to 0$. Note that $\int_\R |f| < \infty$ by integrability of $f$, so we can apply question 1 to get 
\[ \left|F\left(x + \frac{1}{n}\right) - F(x)\right| = \left|\int_\R \mathds{1}_{[x, x + 1/n)} f\right| \leq \int_\R \mathds{1}_{[x, x + 1/n)} |f| \to 0 \]
we can thus take $N > 0$ for which $\left|F\left(x + \frac{1}{N}\right) - F(x)\right| < \varepsilon$. It then follows that for all $h \in (0, 1/N)$, we have 
\[ \left|F\left(x + h\right) - F(x)\right| \leq \int_\R \mathds{1}_{[x, x + h)} |f| \leq \int_\R \mathds{1}_{[x, x + 1/N)} |f| < \varepsilon \]
We can do the same for $A_n = (x - 1/n, x]$ and combine with the above to get a $\delta > 0$ for which $\left|F(y) - F(x)\right| < \varepsilon$ for all $y \in (x - \delta, x + \delta)$. Therefore, $F$ is continuous.
\end{proof}

\section{Royden Chapter 4}

Note that Royden uses the convention that a nonnegative measurable function is integrable when its Lebesgue integral is finite.

\begin{question}[Royden 4.20]
Let $\{f_n\}$ be a sequence of nonnegative measurable functions that converges to $f$ pointwise on $E$. Let $M \geq 0$ be such that 
\[ \int_E f_n \leq M \quad \text{for all } n. \]
Show that $\int_E f \leq M$. Verify that this property is equivalent to the statement of Fatou's Lemma.
\end{question}
\begin{proof}
Since we know Fatou's lemma to be true already, we will simply prove the equivalence of the property to it. If Fatou's lemma holds, then we must have 
\[ \int_E f \leq \liminf_{n \to \infty} \int_E f_n \leq M \]
as required. Conversely, suppose that whenever $f_n \to f$ pointwise on $E$ with $\int_E f_n \leq M$ then $\int_E f \leq M$, and let $f_n \to f$ pointwise on $E$. If $\liminf \int_E f_n$ is infinite, then clearly the statement of Fatou's lemma holds. Otherwise, there is a subsequence $f_{n_k}$ for which $\int_E f_{n_k}$ is a sequence approaching $\liminf\int_E f_n$. Since we still have $f_{n_k} \to f$ pointwise on $E$, for any $\varepsilon > 0$, we can take $K > 0$ such that for all $k \geq K$, we have $\int_E f_{n_k} < \liminf\int_E f_n + \varepsilon$, so by our property, we conclude
\[ \int_E f \leq \liminf_{n \to \infty} \int_E f_n + \varepsilon \]
for all $\varepsilon > 0$. Taking $\varepsilon \to 0$, we prove Fatou's lemma where $f_n \to f$ pointwise everywhere on $E$. Since integrals over sets of measure zero are zero, this clearly extends to convergence a.e.
\end{proof}

\begin{question}[Royden 4.21]
Let the function $f$ be nonnegative and integrable over $E$ and $\epsilon > 0$. Show there is a simple function $\eta$ on $E$ that has finite support, $0 \leq \eta \leq f$ on $E$, and 
\[ \int_E |f - \eta| < \epsilon. \]
If $E$ is a closed, bounded interval, show there is a step function $h$ on $E$ that has finite support and 
\[ \int_E |f - h| < \epsilon. \]
\end{question}
\begin{proof}
By definition of Lebesdue integral, we can take a simple function $\eta \leq f$ such that $\int_E f < \int_E \eta + \varepsilon/2$. This $\eta$ might not have finite support, but we can use it to construct another simple function that has finite support and has integral close to one of $\eta$. Consider the increasing sequence $E_n := E \cap (-n, n)$ and simple functions $\eta_n := \eta\mathds{1}_{E_n}$ with finite support. Since $\bigcup_{n = 1}^\infty E_n = E$, we have $\eta_n \to \eta$ pointwise, and hence by MCT
\[ \int_E \eta = \lim_{n \to \infty} \int_E \eta_n \]
Since the integral is finite, we can take $n$ large enough for which $0 \leq \int_E \eta - \int_E \eta_n < \varepsilon / 2$. By linearity of the integral, we have 
\[ \int_E (f - \eta_n) + \int_E \eta_n = \int_E f \longimplies \int_E |f - \eta_n| = \left(\int_E f - \int_E \eta\right) + \left(\int_E \eta - \int_E \eta_n\right) < \varepsilon \]
as required. 

If $E = [a, b]$ for $a < b$. In this case, any step function on $E$ has finite support, so we don't need to worry about this condition. We first show that the $\eta$ found above can be approximated by a step function:

\begin{lemma}
For any simple function $\varphi$ and $\delta > 0$, we can find a step function $h: E \to [0, \infty]$ and a measurable $F \subseteq E$ such that $\varphi = h$ on $F$ and $m(E \setminus F) < \delta$.
\end{lemma}
\begin{proof}
We note that a sum of finitely many step functions is also a step function, so it suffices to prove this for characteristic functions $\mathds{1}_G$ with $G \subseteq E$ measurable. We take $O = \bigcup_{k = 1}^n I_k$ where $I_k$ are disjoint open intervals such that $m(O \setminus G) + m(G \setminus O) < \delta$ (the existence of such $O$ can be proven by using outer regularity and the fact that open sets are countable unions of open intervals). We then define $h: E \to [0, \infty], x \mapsto \begin{cases}
    1 & x \in O \\ 0 & x \notin O
\end{cases}$, which is a step function because $O$ is a disjoint union of open intervals. We then immediately see that $h(x) = 1 = \mathds{1}_G(x)$ for $x \in G \cap O$ and $h(x) = 0 = \mathds{1}_G(x)$ for $x \in E \setminus (O \cup G)$, so for $F = (G \cap O) \cup (E \setminus (O \cup G)) \subseteq E$, we have $h = \mathds{1}_G$, and 
\begin{align*}
    E \setminus F 
    &= (E \setminus (G \cap O)) \cap (O \cup G) \\
    &= ((E \setminus (G \cap O)) \cap O) \cup ((E \setminus (G \cap O)) \cap G) \\
    &\subseteq ((G \cap O)^c \cap O) \cup ((G \cap O)^c \cap G) \\
    &= (O \setminus G) \cup (G \setminus O) \\
    \longimplies m(E \setminus F) 
    &\leq m(O \setminus G) + m(G \setminus O) < \delta
\end{align*}
as required.
\end{proof}
We take a simple $\eta$ such that $\int_E |f - \eta| < \varepsilon / 2$. Since $\eta$ is simple, it has a finite bound $M > 0$, and we apply our lemma for $\delta = \frac{\varepsilon}{4M}$ to get a step function $h: E \to [0, \infty]$ and a set $F$ such that $m(E \setminus F) < \delta$ and $\eta = h$ on $F$. By the construction in the lemma, we may assume that $h$ is bounded above by $M$, so we have 
\[ \int_E |\eta - h| = \int_{E \setminus F} |\eta - h| \leq 2M \cdot m(E \setminus F) < \frac{\varepsilon}{2} \]
We then have 
\[ \int_E |f - h| \leq \int_E |f - \eta| + \int_E |\eta - h| < \varepsilon \]
\end{proof}

\begin{question}[Royden 4.25]
Let $\{f_n\}$ be a sequence of nonnegative measurable functions on $E$ that converges pointwise on $E$ to $f$. Suppose $f_n \leq f$ on $E$ for each $n$. Show that 
\[
\lim_{n \to \infty} \int_E f_n = \int_E f.
\]
\end{question}

\begin{proof}
\end{proof}

\begin{question}[Royden 4.28]
Let $f$ be integrable over $E$ and $C$ a measurable subset of $E$. Show that 
\[
\int_C f = \int_E f \cdot \chi_C.
\]
\end{question}

\begin{proof}
\end{proof}
    

\end{document}
